% =====================================================================
\section{Introduction}
\label{sec:intro}
% =====================================================================

A gateway that serves large language model (LLM) traffic to many tenants
decides, every few seconds and for every tenant, how many serving replicas
to run, how much semantic-cache memory to grant, and which model tier
should answer the requests the cache misses. In deployed stacks a
different controller owns each decision: the Kubernetes Horizontal Pod
Autoscaler (HPA) or KEDA scales replicas on a utilisation or queue
signal~\cite{k8shpa,keda}, a semantic cache such as GPTCache answers
repeated prompts~\cite{gptcache}, and a router or cascade sends the
remaining requests to a cheaper or a stronger model~\cite{frugalgpt,routellm}.
In the vocabulary of self-adaptive systems, the gateway runs three
feedback loops over one managed system~\cite{kephart2003,weyns2013}. The
loops share capacity: every cache hit removes work from the replicas,
replica congestion slows every miss, and every miss pays its tier's price.
The best setting of each loop in isolation therefore need not be the best
joint setting, and the natural engineering response is one controller
that sets all three knobs together.

Whether that response pays is not settled. Joint configuration of serving
knobs has been built and reported to win: INFaaS selects model variants
together with autoscaling~\cite{infaas}, IPA and Proteus adapt model
variants together with replicas~\cite{ipa,proteus}, Cascadia co-optimises
cascade deployment and routing~\cite{cascadia}, and Zhu et~al.\ design a
cache jointly with a model selector~\cite{zhu2023}. Each compares the joint
design with separately configured or fixed alternatives. Two things are
missing from that evidence. First, controller comparisons are sensitive to
how well the authors tune and instrument their
baselines~\cite{chameleon,zhang2026drl,khilji2026}, and an evaluation that
removes the comparators' handicaps one at a time is rare. Second, few
evaluations isolate what coordination itself is worth, which requires
holding the objective, the forecaster, the actuation limits, the budgets
and the cost model fixed while varying only the search: joint, or per
knob. The question matters to anyone who builds or operates such a
gateway, because a joint controller is costlier to build, to verify and to
debug than three loops that each have a mature implementation.

We built a joint controller, \jcacintro{}, as the planner of \polyforge{},
a Kubernetes-native control plane for multi-tenant LLM gateways, and then
tried to falsify our own hypothesis that joint control is worth building.
Over twelve pre-registered simulation campaigns (33{,}612 runs) we asked
two questions. \emph{RQ1:} does the joint controller beat independently
operated per-layer controllers, each run as a complete policy?
\emph{RQ2:} what does joint search add when everything else is held
fixed? RQ2 is the better isolated of the two; the answer to RQ1 depends on
how the comparators are built and constrained, and we report how it
changed as we corrected them.

The joint controller's advantage did not survive. On RQ2, joint search
lowered the objective by 1.6\% relative to searching each knob with the
other two held fixed: a directional difference, but inside the registered
equivalence margin, and in ten other registered settings the difference
stayed between a 1.4\% loss and a 1.4\% gain. A registered experiment that
varied the actuation step explains why: when one interval may move a
tenant by at most two replicas and one cache level, the per-tenant joint
problem is \emph{locally separable}, so per-knob search finds the joint
move in almost every decision. Coordination grew four- to five-fold with
eight-replica steps, but only on plants where AI work loads the replicas,
and the wide steps made the controller worse. On RQ1, the controller first
appeared to beat HPA by 27.5\% and a layered stack by 60\% of the
objective; both gains came from comparator defects, removed one registered
re-scoring at a time. What then remained rested on the largest cluster, on
actuation lag and on shedding: the controller answered 10.1\% of AI requests
with no model at all, an outage the objective priced lightly.

An audit of all eleven campaigns found leaks on both sides: free outages and
a budget that only the joint controller could keep, which flattered it, and
a comparator running above the cluster caps and an over-punitive never-shed
fallback, which flattered the comparators. A twelfth campaign removed them
together. It charged every refused request, gave a per-layer stack a gateway
budget gate, held every arm to the caps, re-tuned the comparators, and gave
the joint controller its documented fixes, chosen on the comparators' tuning
data. At that level field the joint controller beat no comparator on
synthetic, held-out or lagged demand; it lost to the ungated cheapest-tier
stack, and under actuation lag to the gated stack and to capped HPA, because
the corrected hysteresis that helps it elsewhere pays for every late replica.
It won once: on real BurstGPT demand, among arms that keep the budget, it beat
the budget-gated stack by 2.7\%. With outages priced, coordination rose to
1.3--4.5\% of the objective, and still bought no win.

These are negative results about a controller we designed, and we report
them as the contribution, together with the mechanism that explains them
and the evaluation rules that would have caught our own early errors. The
paper makes six contributions:
\begin{itemize}
\item \textbf{A measurement of what coordination is worth}
  (Sections~\ref{sec:coordination} and~\ref{sec:level},
  Tables~\ref{tab:coord} and~\ref{tab:level}): across eleven registered
  settings joint search changes the objective by between $-1.4\%$ and
  $+1.6\%$ relative to decoupled search, inside the equivalence margin; with
  outages priced it gains 1.3--4.5\%, directional everywhere.
\item \textbf{A final test at a level field}
  (Section~\ref{sec:level}, Table~\ref{tab:level}): with every audited
  leak removed on both sides, the joint controller beats no comparator on
  synthetic, held-out or lagged demand, and beats a budget-gated stack only
  on real BurstGPT demand.
\item \textbf{A mechanism confirmed under pre-registration}
  (Section~\ref{sec:width}, Fig.~\ref{fig:width}): coordination grows with
  how far one interval may move a tenant, and only where the layers share
  capacity; all five registered predictions held.
\item \textbf{An adjudication of the complete-policy comparison}
  (Sections~\ref{sec:chain}--\ref{sec:generalise}, Fig.~\ref{fig:chain}):
  the initial 27.5--60\% gains traced to comparator defects, each removed by
  a registered re-scoring, and what then remained attributed to cluster
  size, actuation lag and shedding; the leaks a final audit found on both
  sides are what the level-field campaign removed.
\item \textbf{An accounting of where the savings come from}
  (Section~\ref{sec:shed}, Table~\ref{tab:served}): request-level counters,
  added without changing any registered metric, show the controller's
  spend advantage is unserved traffic, not cheaper service.
\item \textbf{Evaluation rules and an open artifact}
  (Section~\ref{sec:lessons}, Table~\ref{tab:checklist}): seven rules paired
  with the defects they caught, and code, protocols, run data and scripts
  that regenerate every number in the paper.
\end{itemize}

The simulator's plant is our own model, with serving replicas that are not
GPU model servers, a synthetic cache-hit curve and a synthetic price
ladder; all comparisons are simulated, and a registered live replication
was not funded. Section~\ref{sec:threats} states what these limits leave
open; the results should be read as conditional on that plant.
